\documentclass[11pt]{article}
\usepackage[margin=1in]{geometry}
\usepackage{amsmath, amssymb}
\usepackage{booktabs}
\usepackage{xcolor}
\usepackage{hyperref}
\setlength{\parindent}{0pt}
\setlength{\parskip}{0.8em}

\title{Fate's Edge: Design \& Review Compendium}
\author{Design Team}
\date{}

\begin{document}

\maketitle

\tableofcontents
\newpage

\section*{Introduction}
This document compiles three essential resources for designers and reviewers evaluating \textit{Fate's Edge}:
\begin{enumerate}
    \item \textbf{Reviewer's Guide} – A framework for assessing the game's narrative-first design, progressive complexity, and storytelling partnership model.
    \item \textbf{Mechanics Breakdown} – A detailed analysis of the core resolution engine, resource economy, and tactical levers.
    \item \textbf{Tiered Heuristics} – Optimal play strategies for Tiers I–IV, grounded in probability and resource scarcity.
\end{enumerate}

All content is presented as raw LaTeX for integration into production documents. No editorial commentary is included; the material is forward-facing for fellow designers and reviewers.

\newpage

\part{Reviewer's Guide: Understanding Fate's Edge}

\textit{This section will be removed in production.}

\section{Reframing the Evaluation}

Fate's Edge is not a traditional tabletop RPG to be evaluated by comprehensive rule mastery. Instead, it is a \textbf{narrative partnership} designed to evolve with your table's storytelling sophistication. Approach this review through the lens of \textbf{progressive storytelling enhancement} rather than static mechanical complexity.

\section{Key Evaluation Framework}

\subsection{Start Simple, Discover Naturally}
\begin{itemize}
    \item \textbf{Session 1}: Core mechanic + 1–2 regional generators = complete playable experience
    \item \textbf{Session 5}: Comfortable integration of multiple subsystems
    \item \textbf{Session 20}: Intuitive mastery where mechanics become invisible
\end{itemize}

\subsection{The Partnership Model}
Rather than ``learning a game,'' players develop a ``storytelling partnership'' where:
\begin{itemize}
    \item Rules serve narrative consequences, not simulation accuracy
    \item Complexity emerges from interaction, not individual component sophistication
    \item Prep burden shifts from authoring content to facilitating discovery
\end{itemize}

\section{Common Misunderstandings to Address}

\subsection{Content Volume vs. Discovery Process}
\textbf{Misconception}: ``Too much content to master'' \\
\textbf{Reality}: Content is designed for gradual discovery through play

\begin{quote}
\textit{Reviewer Guidance:} ``Imagine a library where you only need to know about the books relevant to today's story. The rest exist to be discovered when narrative ambition demands them.''
\end{quote}

\subsection{Prep-Light vs. Prep-Free}
\textbf{Misconception}: ``No prep means unstructured play'' \\
\textbf{Reality}: Prep is embedded in the constraint lattice

\begin{quote}
\textit{Reviewer Guidance:} ``The system replaces authored prep with guided storytelling. Every card draw is improvisation with guardrails.''
\end{quote}

\subsection{Mechanical Simplicity vs. Narrative Sophistication}
\textbf{Misconception}: ``Simple mechanics = shallow stories'' \\
\textbf{Reality}: Sophistication emerges from constraint interaction

\begin{quote}
\textit{Reviewer Guidance:} ``Like jazz with a vast repertoire, the same simple rules that handle Session 1 magic effortlessly scale to manage Session 50 world-changing consequences.''
\end{quote}

\section{Evaluation Methodology}

\subsection{Progressive Complexity Assessment}
\begin{enumerate}
    \item \textbf{Tier I Play} (30–40 XP): Evaluate session-to-session narrative flow with minimal subsystems
    \item \textbf{Tier III Play} (90–150 XP): Assess how character growth naturally expands storytelling scope
    \item \textbf{Tier V Play} (220+ XP): Examine how simple rules handle complex, multi-threaded narratives
\end{enumerate}

\subsection{Mastery Curve Analysis}
\textbf{Key Metric}: Mastery effort should decrease as narrative complexity increases

\textbf{Success Indicators}:
\begin{itemize}
    \item Session 1: Clear, engaging storytelling with minimal reference
    \item Session 10: Comfortable integration of multiple subsystems
    \item Session 30: Mechanics become invisible; focus entirely on narrative
\end{itemize}

\section{What to Look For}

\subsection{Narrative-First Design}
\begin{itemize}
    \item Every mechanical element should serve story consequences
    \item Player agency should manifest through narrative choices, not optimization
    \item GM authority should prioritize ``what happens'' over ``what are the rules''
\end{itemize}

\subsection{Constraint-Based Creativity}
\begin{itemize}
    \item Random elements should create coherent, not chaotic, storytelling
    \item Regional generators should maintain thematic consistency while enabling variety
    \item Mechanical constraints should inspire creativity, not limit options
\end{itemize}

\subsection{Scalable Sophistication}
\begin{itemize}
    \item Same core mechanics should handle intimate scenes and epic narratives
    \item Character progression should expand narrative possibilities, not mechanical options
    \item World complexity should emerge from player actions, not pre-authored content
\end{itemize}

\section{Red Flags vs. Design Features}

\subsection{Actual Red Flags}
\begin{itemize}
    \item Mechanics that require constant reference during play
    \item Subsystems that don't integrate naturally with core philosophy
    \item Progression that adds mechanical burden rather than narrative freedom
\end{itemize}

\subsection{Misinterpreted Design Features}
\begin{itemize}
    \item \textbf{High content volume} = Creative resource library, not mastery burden
    \item \textbf{Multiple magic paths} = Narrative vocabulary expansion, not mechanical complexity
    \item \textbf{Regional generators} = Inspiration engines, not reference requirements
\end{itemize}

\section{Recommended Review Approach}

\subsection{Session 1–3 Evaluation}
Focus on:
\begin{itemize}
    \item How quickly engaging stories emerge
    \item Clarity of core mechanical principles
    \item Integration of generated content into coherent narratives
\end{itemize}

\subsection{Session 10+ Evaluation}
Assess:
\begin{itemize}
    \item How naturally complexity emerges from player choices
    \item Integration of multiple subsystems without mechanical overhead
    \item Evolution from ``using rules'' to ``telling stories''
\end{itemize}

\subsection{Long-term Assessment}
Consider:
\begin{itemize}
    \item Whether the system grows with narrative ambition
    \item If mastery effort genuinely decreases over time
    \item How world complexity emerges organically from play
\end{itemize}

\section{The Core Truth}

Fate's Edge succeeds not by being simple, but by being \textbf{designed for simplicity to emerge naturally from complex storytelling}. Evaluate it not as a comprehensive system to master, but as a storytelling partnership that grows in sophistication at exactly the same rate your table develops narrative ambition.

\begin{quote}
\textbf{This isn't a game you learn – it's a storytelling evolution you grow into.}
\end{quote}

\newpage

\part{Mechanics Breakdown}

\section{The Core Philosophy: Beyond Success/Failure}

Most RPGs reduce outcomes to a binary: did you succeed or fail? Fate's Edge rejects this. Instead, it asks: \textit{What happens next?} This single question reshapes every mechanic.

The design mantra: \textbf{Every roll must move the story forward—no matter the result.} Even a ``failure'' creates new opportunities, while success often introduces complications. This eliminates dead ends and ensures the narrative constantly evolves.

\textbf{Why this matters to strategists}: The system transforms dice from pass/fail meters into narrative catalysts. Your goal isn't to maximize success rate—it's to manage the flow of consequences so that both triumphs and setbacks drive meaningful story.

\section{The Dice Engine: d10s with Purpose}

\subsection{How Rolls Work}
\begin{itemize}
    \item Roll a pool of d10s: \textbf{Attribute (1–5) + Skill (0–5)}
    \item Each die showing 6–10 = 1 success; a natural 10 = 2 successes
    \item Target number: Difficulty Value (DV) from 2 (routine) to 5+ (extreme). Default DV = 3
\end{itemize}

\subsection{Why d10s?}
\begin{itemize}
    \item Clean 50\% success chance per die (6–10 = 5 out of 10 faces)
    \item Enables quick mental math: 4 dice → 50\% chance of at least one 10
    \item The ``10 = 2 successes'' mechanic creates meaningful variance without extra rolls
\end{itemize}

\subsection{Key Probability Insight}
With 4 dice (a common sweet spot):
\begin{itemize}
    \item Chance of at least one success: 93.75\% (only 6.25\% miss rate)
    \item Chance of at least one 10: 34.4\% (triggers bonus effects)
    \item Chance of rolling a 1 (for GM): 34.4\% (gives the GM narrative leverage)
\end{itemize}

\textbf{Strategic implication}: The first points in attributes/skills give huge returns (going from 2→3 dice cuts miss rate from 25\% → 12.5\%). Beyond 4 dice, diminishing returns kick in sharply—miss rate drops from 6.25\% → 3.9\% for 5 dice, but GM leverage increases significantly. Smart players cap raw dice at 4 and use positioning to boost effective pool size.

\section{Position: The Tactical Lever}

Position (Dominant/Controlled/Desperate) doesn't just modify difficulty—it reshapes your probability distribution before you roll:

\begin{table}[h]
\centering
\begin{tabular}{lll}
\toprule
\textbf{Position} & \textbf{Narrative Frame} & \textbf{Mechanical Effect} \\
\midrule
Dominant & You press your advantage & Re-roll one failure \\
Controlled & Balanced, normal risk & No re-rolls \\
Desperate & You act under duress & Re-roll one success \\
\bottomrule
\end{tabular}
\end{table}

\subsection{What this does mathematically}
\begin{itemize}
    \item \textbf{Dominant}: Cuts the left tail of your success distribution. Reduces chance of zero successes dramatically (e.g., 25\% → 6.25\% at 2 dice).
    \item \textbf{Desperate}: Increases volatility – more likely to get extreme results (0 or 2+ successes). Favors high-skill players who can turn a re-rolled 6 into a 10.
\end{itemize}

\textbf{Critical nuance}: Position adjustments stack as dice. If already Dominant, further improvements become +1 die (not additional re-rolls). This prevents abuse while keeping positioning meaningful.

Why not adjust DV? Because position reflects narrative risk, not raw task difficulty. Picking a lock while guards approach (Desperate) feels different than in a safe room (Dominant)—even if the DV is identical. This forces players to engage with fiction tactically.

\section{The Outcome Matrix: Four Paths Forward}

Every roll lands in one of four states—all create narrative momentum:

\begin{table}[h]
\centering
\begin{tabular}{ll}
\toprule
\textbf{Roll Result} & \textbf{Outcome} \\
\midrule
Successes ≥ DV, SB = 0 & Clean Success. Achieve intent without complication. \\
Successes ≥ DV, SB > 0 & Success with SB. Achieve intent; GM spends SB for a twist. \\
0 < Successes < DV & Partial. Make progress; player gains 1 Boon. \\
Successes = 0 & Miss. No progress; player gains 2 Boons; GM escalates with SB. \\
\bottomrule
\end{tabular}
\end{table}

Why this eliminates dead ends:
\begin{itemize}
    \item A Partial isn't failure—it's meaningful progress (e.g., ``The lock is stiff but you think you can get it if you keep trying'').
    \item A Miss generates resources (Boons) that let you influence the next roll.
    \item The GM must spend SB to complicate the situation—nothing is ever ``nothing happens.''
\end{itemize}

\textbf{Strategic depth}: The system creates a resource loop:
\begin{itemize}
    \item Boons (from Partial/Miss) let you re-roll dice or improve position
    \item Story Beats (SB) (from 1s rolled) let the GM introduce complications
    \item Smart players embrace small failures early to stockpile Boons for critical moments—but Boons cap at 5 and reset partially between scenes, preventing hoarding.
\end{itemize}

\section{Resources: Boons (Player) and Story Beats (GM)}

\subsection{Boons: Your Tactical Reserve}
Boons represent momentum, luck, and narrative leverage. Spend them to:
\begin{itemize}
    \item Re-roll a single die (1 Boon)
    \item Improve position by one step before rolling (1 Boon)
    \item Activate an off-screen Asset (1 Boon)
    \item Convert 2 Boons → 1 XP (once per session, max 2 XP/session)
\end{itemize}

\textbf{The XP conversion trap}: Converting Boons to XP is rarely optimal. One Boon via re-roll gives 0.4 expected successes—often worth more than the XP gained from those successes. Only convert when:
\begin{itemize}
    \item You're at the 5-Boon cap (about to lose excess)
    \item The session is ending (Boons reset to 2)
\end{itemize}

\textbf{Optimal use}: Always keep 1 Boon for critical re-rolls. Spending 1 Boon to improve position (e.g., Controlled → Dominant) before a big roll is often better than re-rolling after—it affects your entire pool.

\subsection{Story Beats: GM's Narrative Fuel}
Every 1 rolled gives the GM 1 Story Beat. Spend SB to:
\begin{itemize}
    \item \textbf{1 SB}: Minor complication (noise, distraction, +1 clock segment)
    \item \textbf{2 SB}: Moderate complication (alarm raised, lose advantage, worsen position)
    \item \textbf{3+ SB}: Major complication (reinforcements, scene shift, broken asset)
\end{itemize}

\textbf{The high-level shift}: At lower tiers, the GM spends SB sparingly (1–2/scene). At tiers III+ (90+ XP), players routinely roll 6–8 dice—generating 0.47–0.57 expected SB per roll. In a 4-player round, that's 2–2.3 SB generated. Smart GMs spend this aggressively:

\begin{quote}
``You pick the lock cleanly—but the noise alerted the patrol. Reinforcement [6] tick +2.''
``You disarm the trap—but the recoil jars your aim. Your next shot is at –1 die.''
\end{quote}

Result: Even successes create meaningful complications. Players learn: Success means the world reacted—not that you're safe.

\section{Harm and Fatigue: The Anti-Death Spiral}

Most games punish injury with worsening rolls—a death spiral. Fate's Edge does the opposite:
\begin{itemize}
    \item \textbf{Fatigue}: Tracks exertion (max = Body). Each point worsens position (Dominant → Desperate). If already Desperate, each fatigue = –1 die.
    \item \textbf{Harm}: Injury severity (1 = –1 die on related actions, 2 = –2 dice on most actions, 3 = incapacitated).
\end{itemize}

The key innovation: \textbf{Taking harm clears all fatigue.} The shock of a wound resets exhaustion—creating dramatic swings (e.g., a fatigued character who takes harm suddenly finds second wind).

\subsection{Armor's Role}
Armor converts harm to fatigue:

\begin{table}[h]
\centering
\begin{tabular}{lccc}
\toprule
\textbf{Armor Type} & \textbf{Harm 1} & \textbf{Harm 2} & \textbf{Harm 3} \\
\midrule
Light  & Fatigue 1 & Harm 1 + Fatigue 1 & Harm 2 + Fatigue 1 \\
Medium & Fatigue 1 & Fatigue 1 & Harm 2 + Fatigue 1 \\
Heavy  & Fatigue 1 & Fatigue 1 & Fatigue 2 \\
\bottomrule
\end{tabular}
\end{table}

\textbf{Strategic implication}: Heavy armor turns lethal blows into manageable exhaustion—but is less efficient against multiple small hits. Smart players:
\begin{itemize}
    \item Use heavy armor against expected high-damage threats
    \item Accept light/medium armor for versatile threat environments
    \item Sometimes seek low-risk harm (e.g., trapping a hand in a door) to reset fatigue when overwhelmed
\end{itemize}

\section{Clocks: Visible Tension Tracking}

Clocks track progress toward an event (4–10 segments). They advance via:
\begin{itemize}
    \item Misses/Partials (often +1 segment)
    \item GM spending SB (usually 1 segment per SB)
    \item Player actions (e.g., successful skill checks)
\end{itemize}

Why they matter strategically: Clocks turn abstract threats into tangible resources. A ``Barrow Collapse [6]'' clock means:
\begin{itemize}
    \item You have 6 ``safe'' actions before the entrance seals
    \item Each failed roll or GM SB spend brings you closer to disaster
    \item Smart players treat clocks like countdown timers—allocating actions to delay or accelerate based on goals
\end{itemize}

\textbf{High-level insight}: At higher tiers, clocks become the primary threat meter. The GM doesn't just add complications—they use clocks to make pressure visible and inevitable. Players must manage clock state alongside their resources.

\section{Magic: The TAGS System}

Free casting uses TAGS (keywords) to build spells. DV is calculated simply:
\[
\text{DV} = 1 + \text{number of TAGS}
\]
with dangerous TAGS (Leap, Fold, Gate, Transform, Dominate) adding +2 DV each.

\subsection{Example Spells}
\begin{table}[h]
\centering
\begin{tabular}{lll}
\toprule
\textbf{Spell} & \textbf{TAGS} & \textbf{DV} \\
\midrule
Light & [Illuminate] & 2 \\
Heal & [Heal] & 2 \\
Shield of Air & [Air][Defend] & 3 \\
Fire Wall & [Fire][Wall] & 3 \\
Teleport & [Fold] (dangerous) & 1+2=3 \\
\bottomrule
\end{tabular}
\end{table}

Critical clarification:
\begin{itemize}
    \item \texttt{[Air][Defend]}: Pure elemental build (DV 3)
    \item \texttt{[Ward]}: Uses ward mechanics (DV = [Ward]'s base, usually 3) creates an integrity clock—the shield can take hits before failing. This is superior for sustained defense—\texttt{[Air][Defend]} is one-time only.
\end{itemize}

\subsection{Why this system wins}
\begin{itemize}
    \item Transforms spell design into a cost-benefit equation
    \item Rewards creativity within clear constraints (e.g., ``Firebolt'' = [Fire] DV 2; add [Ranged] for distance = DV 3)
    \item Scales risk: More TAGS = higher DV = more chance of backlash
    \item Backlash is thematic (fire backlash = smoke → blaze; earth = dust → fissure)
\end{itemize}

\textbf{Smart caster play}: Build minimum TAGS for effect, then:
\begin{itemize}
    \item Use positioning to improve effective DV
    \item Accept manageable backlash (e.g., fatigue 1) for high-impact spells
    \item Save high-risk spells for moments when you can mitigate consequences
\end{itemize}

\section{Advancement: The XP Economy}

XP costs scale to encourage meaningful choices:

\begin{table}[h]
\centering
\begin{tabular}{lll}
\toprule
\textbf{Improvement} & \textbf{Cost} & \textbf{Downtime} \\
\midrule
Attribute increase & new rating × 3 XP & new rating (days) \\
Skill increase & new level × 2 XP & new level (days) \\
Minor Asset (tool, contact) & 4 XP & 1 day \\
Standard Asset (safehouse, network) & 8 XP & 1 week \\
Major Asset (fortress, guildhall) & 12 XP & 1 month \\
Minor Talent (edge case) & 2–4 XP & 3–10 days \\
Major Talent (core ability) & 5–8 XP & 1–2 weeks \\
\bottomrule
\end{tabular}
\end{table}

\subsection{Design rationale}
\begin{itemize}
    \item Attributes cost more because they affect multiple skills (e.g., high Body aids melee, athletics, stealth)
    \item Skills cost less because they're narrower in application
    \item Creates natural progression: Broad attribute foundation → Skill specialization → Talent mastery
\end{itemize}

\subsection{Diminishing returns in action}
\begin{itemize}
    \item Raising an attribute from 4→5: 15 XP for +1 die on 3–5 skills
    \item Raising a skill from 4→5: 10 XP for +1 die on only that skill
    \item Optimal path: Max key attributes to 4 first (60 XP for all four at 4), then specialize skills to 3–4, then take talents
\end{itemize}

\subsection{Tiers of play (XP ranges)}
\begin{itemize}
    \item \textbf{Novice (0–40)}: Local reputation, prestige locked
    \item \textbf{Seasoned (41–90)}: Regional notice, prestige abilities unlock
    \item \textbf{Veteran (91–150)}: National influence, second follower slot suggested
    \item \textbf{Paragon (151–220)}: Movers and shakers, rivals emerge
    \item \textbf{Mythic (221+)}: Legendary status, kingdoms respond
\end{itemize}

At Veteran+: You'll regularly roll 8+ dice. The game responds not by inflating DV, but by making the GM's SB spend more impactful—reinforcements arrive faster, complications cascade, and success requires meaningful sacrifice.

\section{Design Philosophy Summary}

\begin{table}[h]
\centering
\begin{tabular}{ll}
\toprule
\textbf{Goal} & \textbf{Mechanic} \\
\midrule
No wasted rolls & 4-state outcome matrix (always creates narrative) \\
Failure drives story & Miss generates Boons; GM spends SB \\
Risk vs. reward & Larger dice pools = more success and more 1s (GM leverage) \\
Tactical positioning & Position reshapes probability distribution pre-roll \\
No death spiral & Harm clears fatigue \\
Visible tension & Clocks track progress toward events \\
Creative magic & TAGS system: DV = 1 + TAGS (dangerous +2) \\
Balanced advancement & Attribute cost vs. skill cost; clear XP tiers \\
\bottomrule
\end{tabular}
\end{table}

\begin{quote}
\textbf{Final strategist's tip}: The most powerful upgrade isn't more dice—it's effective dice via positioning. Spending 1 Boon to improve position (Controlled → Dominant) before a roll is equivalent to +2 dice in miss reduction—but costs 0 XP. Save your XP for attributes that cover multiple skills, then use positioning and Boons to hit the 4–5 effective dice sweet spot where miss rates plummet but GM leverage stays manageable.
\end{quote}

\newpage

\part{Tiered Heuristics}

\section{Introduction to Tiered Play}

The following heuristics are derived from probability simulations and extensive playtesting. They provide optimal strategies for each tier of play (I–IV), grounded in the game's core design principle: \textbf{DV does not inflate; instead, SB and Boon scarcity drive difficulty.}

\subsection{Common Assumptions}
\begin{itemize}
    \item Default DV = 3 (routine) to 4 (challenging)
    \item Position starts Controlled unless fiction dictates otherwise
    \item Players use bonds and assets as intended
    \item GM spends SB aggressively, following the SB spend tables
\end{itemize}

\section{Tier I (0–40 XP) Optimal Heuristics}

\subsection{Combat – Melee vs. Ranged}
\textbf{Scenario}: 1v1 against a bandit (Cap 2, Defense DV 3, Harm 1 per hit). Tier I character has Body 3 + Melee 2 = 5 dice.

\begin{table}[h]
\centering
\begin{tabular}{lccc}
\toprule
\textbf{Dice} & \textbf{P(success ≥3)} & \textbf{P(Miss 0)} & \textbf{Expected SB} \\
\midrule
5 & ~72\% & ~3\% & 0.5 \\
\bottomrule
\end{tabular}
\end{table}

\textbf{Heuristic}: 5 dice gives 72\% success rate. For high reliability (≥85\%), need 6 dice (Body 3 + Melee 3).

\textbf{Optimal Tier I melee build}: Body 3, Melee 3 (6 dice). Cost: Body 3 (9 XP) + Melee 3 (6 XP) = 15 XP. Leaves room for other skills.

\textbf{Position impact}: Dominant (re-roll failures) effectively adds ~1.5 dice. Desperate hurts reliable characters more than it helps – avoid unless you have Boons to spare.

\textbf{Boons}: One Boon spent to re-roll a failure on a 5-die pool increases success rate from 72\% → ~82\% (worth ~10\% absolute). Best used when you have at least one success and need one more.

\textbf{Takeaway}: At Tier I, prioritize getting your primary combat skill to 3 (6 XP) and attribute to 3 (9 XP). Don't spread thin.

\subsection{Social – Persuasion}
\textbf{Scenario}: Convince a guard to let you pass (DV 3). Presence 2 + Sway 2 = 4 dice.

\begin{table}[h]
\centering
\begin{tabular}{lccc}
\toprule
\textbf{Dice} & \textbf{P(success ≥3)} & \textbf{P(Partial)} & \textbf{P(Miss)} \\
\midrule
4 & ~54\% & ~38\% & ~8\% \\
\bottomrule
\end{tabular}
\end{table}

\textbf{Heuristic}: 4 dice is unreliable. Spend a Boon to improve Position (Controlled → Dominant) before rolling: re-roll failures effectively adds ~1 die, raising success to ~66\%. Better than re-rolling after.

\textbf{Optimal social build}: Presence 3 (9 XP) + Sway 3 (6 XP) = 15 XP. 6 dice gives ~87\% success. But Tier I resources are tight. For secondary role, Presence 2 + Sway 2 (4 dice) + Boon for Dominant is acceptable.

\subsection{Stealth}
\textbf{Scenario}: Sneak past a sentry. DV 3, Wits 3 + Stealth 2 = 5 dice.

\textbf{Heuristic}: 5 dice at Controlled gives ~72\% success. Sentries often have Notice pools of 3–4 dice; opposed roll? Fate's Edge uses DV set by GM, not opposed. So 5 dice is solid.

\textbf{Optimal}: Wits 3 (9 XP) + Stealth 3 (6 XP) = 15 XP gives 6 dice (~87\%). Overkill? Not if you're the scout.

\textbf{Position tricks}: Dominant is huge for stealth – one Miss can alert the whole camp. Spend 1 Boon before rolling to upgrade position.

\subsection{Investigation / Lore}
\textbf{Scenario}: Research a forgotten ritual. DV 3, Wits 3 + Lore 2 = 5 dice.

\textbf{Heuristic}: Same as stealth. But note: failures here generate SB that GM can use for narrative twists (e.g., you find the info but also alert a guardian). Acceptable.

\textbf{Optimal}: If this is your niche, Lore 3 (6 XP) is cheap. Wits 3 + Lore 3 = 6 dice.

\subsection{Travel / Survival}
\textbf{Scenario}: Navigate a broken road in Acasia. DV 3, Wits 2 + Survival 2 = 4 dice.

\textbf{Heuristic}: 4 dice → 54\% success. Travel often allows assistance (party member helps → +1 die up to +3 max). With 2 helpers, you can reach 6 dice (87\%). Don't invest heavily in Survival unless you're the guide.

\subsection{Resource Management – Boons}
\begin{itemize}
    \item You gain 1 Boon on a Partial, 2 Boons on a Miss. Hold up to 5 Boons, reset to 2 at scene end.
    \item \textbf{Optimal Boon use}: 
        \begin{enumerate}
            \item \textbf{Before rolling}: Spend 1 Boon to upgrade Position (Controlled → Dominant) – best value.
            \item \textbf{After rolling}: Reroll a failure – second best.
            \item Activate Asset – situational.
            \item Convert to XP – only at session end if you're at cap.
        \end{enumerate}
    \item \textbf{Heuristic}: Never let a Partial's Boon go to waste – use it on the next roll to improve position. Cycle: Partial → gain 1 Boon → spend it on next roll for Dominant → higher success rate → fewer Misses.
    \item Boons from Misses (2) are painful, but if you have Boons stored, you can recover. Misses are more acceptable when you have 2+ Boons already.
\end{itemize}

\subsection{Damage Output vs. Survival}
\textbf{Tier I character}: 3 Harm capacity (Harm 3 = incapacitated). Body 3 gives Fatigue track size 3.

Armor choice (from earlier table): Light armor fine if high Body (Fatigue track 4+). Medium best for all-rounders. Heavy for tanks.

Damage output: At Tier I, max damage is ~2 Harm per hit (with a 10). To down a Tier I foe (3 Harm) in 1–2 hits, need +1 Effect or critical hits.

\textbf{Takeaway}: +1 Effect (from talents or positioning) is more valuable than +1 die for damage. A shove to make prone (gives +1 Effect) is often better than attacking twice.

\subsection{Clock Management}
\textbf{Scenario}: 6-segment ``Reinforcements Arrive'' clock. Each failed roll ticks +1.

\textbf{Heuristic}: Treat clocks as resources. If you have high dice pools (≥6), you can risk ticking the clock because you'll likely succeed. If you have low pools (≤4), avoid rolls that advance clocks – use alternate approaches (stealth, negotiation) that don't trigger the clock.

Clock size determination (from generators): 2–5 rank = 4 segments; 6–10 = 6; J/Q/K = 8; A = 10. At Tier I, 6-segment clocks are common.

\subsection{Skill Synergies}
At Tier I, focus on 2–3 key skills at level 2–3. Optimal spread for 32–36 XP build:

\begin{table}[h]
\centering
\begin{tabular}{llll}
\toprule
\textbf{Role} & \textbf{Primary Attribute} & \textbf{Key Skills (lvl 2-3)} & \textbf{XP (att+skills)} \\
\midrule
Fighter & Body 3 & Melee 3, Athletics 2, Command 1 & 21 \\
Face & Presence 3 & Sway 3, Insight 2, Deception 1 & 21 \\
Scout & Wits 3 & Stealth 3, Survival 2, Subterfuge 1 & 21 \\
Scholar & Wits 3 & Lore 3, Arcana 2, Investigation 1 & 21 \\
Healer & Spirit 3 & Medicine 3, Survival 2, Insight 1 & 21 \\
\bottomrule
\end{tabular}
\end{table}

21 XP on core attributes+skills leaves 11–15 XP for a minor talent (2–4 XP), a minor asset (4 XP), and/or a second skill at 2 (4 XP). This is the sweet spot.

\subsection{Optimal Starting Build Template (Tier I, 32–36 XP)}
\begin{verbatim}
Attributes: Primary 3, Secondary 2, Tertiary 2, Dump 1 (cost: 9+6+6+3 = 24 XP)
Skills: Primary skill 3 (6 XP), Secondary skill 2 (4 XP), Tertiary skill 1 (2 XP) = 12 XP
Total: 36 XP – exactly at max.
\end{verbatim}
If you want an asset or talent, drop a skill level.

\subsection{Tier I Heuristics Summary}
\begin{enumerate}
    \item Target 6 dice for primary action (Attribute 3 + Skill 3) → ~87\% success at DV 3.
    \item 5 dice (3+2) is acceptable (72\%) – use Boons to upgrade position.
    \item Spend Boons before rolling to upgrade Position – best ROI.
    \item Don't hoard Boons – you reset to 2 at scene end.
    \item Embrace Partials – they give a Boon and keep the story moving.
    \item Misses are acceptable when you have Boons to buffer, but avoid on critical rolls.
    \item Use assistance – helping adds up to +3 dice.
    \item Clocks are timers – track visibly. Use high-dice rolls to advance safely.
    \item Armor choice: Medium is best for most Tier I characters.
    \item Skill spread: 3/2/1 optimal for 36 XP. For an asset, drop to 3/1/1 or 2/2/1.
\end{enumerate}

\section{Tier II (41–90 XP) Optimal Heuristics}

\subsection{Combat}
\textbf{Scenario}: Fighting a Cap 3 elite (Defense DV 4, Harm 2 per hit). Tier II baseline: Body 3 + Melee 3 = 6 dice, often boosted to 7–8 via talents/assets.

\begin{table}[h]
\centering
\begin{tabular}{lcccc}
\toprule
\textbf{Dice} & \textbf{P(success ≥4)} & \textbf{P(Partial)} & \textbf{P(Miss)} & \textbf{Exp SB} \\
\midrule
6 & ~59\% & ~38\% & ~3\% & 0.6 \\
7 & ~68\% & ~30\% & ~2\% & 0.7 \\
8 & ~76\% & ~23\% & ~1\% & 0.8 \\
\bottomrule
\end{tabular}
\end{table}

\textbf{Heuristic}: Want at least 7 dice for ~68\% success against DV 4. 8 dice (76\%) is reliable.

\textbf{Optimal Tier II melee}: Body 4 (12 XP from 3→4) + Melee 4 (8 XP from 3→4) = 20 XP upgrade from Tier I. Total dice = 8.

\textbf{Position impact}: Dominant on 8 dice pushes success to ~87\% against DV 4. Desperate risky.

\textbf{Boons}: One Boon for Dominant before rolling still optimal. Rerolling single failure on 8 dice less efficient (76\% → ~82\%).

\textbf{Talent synergy}: ``Battle Instincts'' (reroll failed defense) huge. ``Exceptional Coordination'' (follower gives +4 dice) pushes to 12 dice (overkill but fun).

\subsection{Ranged Combat}
Similar dice math. Ranged often allows Controlled/Dominant from high ground. Ammo management via Supply clock. Minor asset ``Quiver of Endless Arrows'' (4 XP) worthwhile.

\subsection{Social – High-Stakes Negotiation}
\textbf{Scenario}: DV 4, Presence 3 + Sway 3 = 6 dice → 59\% success. Too low. Upgrade to Presence 4 + Sway 3 (7 dice → 68\%) or Presence 3 + Sway 4 (7 dice). Better: Presence 4 + Sway 4 (8 dice → 76\%).

Optimal social: Presence 4 (12 XP) + Sway 4 (8 XP) = 20 XP. Add ``Silver Tongue'' (4 XP) for +1 die → 9 dice (~83\%). Social assets: ``Noble Contact'' (Minor, 4 XP) or ``Diplomatic Papers'' grant +1 Position.

\subsection{Stealth – Infiltration}
\textbf{Scenario}: Sneak past magical ward (DV 4). Wits 4 + Stealth 3 = 7 dice → 68\%. Acceptable. For DV 5 (magical ward + sentry), need 8–9 dice.

Optimal stealth: Wits 4 + Stealth 4 (8 dice). ``Shadow's Cloak'' asset (advantage in dim light) adds effective +2 dice. Talent ``Shadow Dance'' (10 XP) lets re-enter stealth after kill – very strong.

\subsection{Investigation / Lore}
\textbf{Scenario}: DV 4, Wits 4 + Lore 3 = 7 dice (68\%). Often take extra time (upgrade position) or use ``Library'' asset (Minor, +1 die). Optimal: Wits 4 + Lore 4 = 8 dice (76\%). ``Lorekeeper'' talent (4 XP) invaluable for plot-critical moments.

\subsection{Travel / Survival}
\textbf{Scenario}: DV 4, Wits 3 + Survival 3 = 6 dice → 59\%. But group assistance can push to 9 dice (87\%). Survival 3 sufficient if party cooperates. Use Guide's Braid or similar assets to reduce DV.

\subsection{Resource Management – Boons}
At Tier II, larger dice pools roll more 1s: 0.6–0.8 SB per roll. GM will spend aggressively. Boon sources: Partials (main) and bonds. Misses rare.

\textbf{Heuristic}: Boon scarcity increases. Use Boons for Dominant position rather than re-rolls. Attempt lower-stakes roll first to generate a Boon (Partial), then use that Boon to upgrade position on critical roll.

\subsection{Damage Output vs. Survival}
Tier II foes: Harm 3–4, may have Armor 1. Player Harm capacity: Body 3 → Fatigue 3. With medium armor, convert Harm 2 → Fatigue 1.

Damage output: 8 dice + good weapon (+1 Harm from talent) → expected ~2–3 Harm per hit. Down Tier II foe in 1–2 hits.

Optimal damage talents: ``Backstab'' (8 XP), ``Radiant Smite'', ``Blood Frenzy''. Defensive: ``Battle Instincts'', ``Unnatural Resilience''.

\subsection{Clock Management}
Clocks often 8 segments. Use high-dice characters to advance primary clock, lower-dice to manipulate secondary clocks via non-dice actions (bribes, distractions).

\subsection{Skill Synergies \& Build Templates (70–80 XP)}
XP budget from Tier I (36 XP) plus 34–44 XP earned = 70–80 XP.

\textbf{Upgrade priorities}:
\begin{enumerate}
    \item Raise primary attribute to 4 (12 XP)
    \item Raise primary skill to 4 (8 XP)
    \item Raise secondary attribute to 3 (9 XP)
    \item Major talent (6–10 XP)
    \item Minor asset (4 XP) or Cap 3 follower (9 XP)
\end{enumerate}

\textbf{Example Fighter (80 XP)}: From Tier I baseline, upgrade Body 3→4 (12), Melee 3→4 (8), Wits 2→3 (9). Add ``Battle Instincts'' (6) and minor asset (4) = 39 XP from Tier I. Left 41 XP for Presence, follower, etc.

\textbf{Example Face (80 XP)}: Presence 3→4 (12), Sway 3→4 (8), Wits 2→3 (9). Add ``Silver Tongue'' (4) and ``Diplomatic Papers'' (4) = 37 XP from Tier I. Left 43 XP.

\subsection{Tier II Heuristics Summary}
\begin{enumerate}
    \item Target 8 dice for primary actions (Attribute 4 + Skill 4) → 76\% vs DV 4.
    \item Boons scarcer (Misses rare). Use to upgrade Position, not re-rolls.
    \item Partials are main Boon source – a Partial is success with cost; take the Boon and use next roll.
    \item SB generation increases (0.8–1.0 per roll). Expect regular complications.
    \item Get a major talent (6–10 XP) – affordable and defining.
    \item Raise secondary attribute to 3 – broadens versatility.
    \item Assets matter – Minor asset (4 XP) gives free off-screen effect or +1 die situationally.
    \item Followers (Cap 3, 9 XP) add +3 assist dice – very strong.
    \item Clocks often 8 segments – plan 8–10 actions. High-dice advance, low-dice manipulate.
    \item Defense matters – enemies hit harder. Take defensive talent or invest in medium/heavy armor.
\end{enumerate}

\section{Tier III (91–150 XP) Optimal Heuristics}

\textbf{Core GM Principle}: Do not artificially inflate DV. Keep standard DV at 3. Instead, spend the glut of SB (from larger dice pools) on more pressure, interesting twists, and cascading complications.

\textbf{Consequence}: Players succeed more often, but every success comes with a cost. Boons become scarce because Misses are rare. Characters are more powerful (more dice) but more brittle (less Boon cushion).

\subsection{The Probability Shift}
\begin{table}[h]
\centering
\begin{tabular}{lcccc}
\toprule
\textbf{Dice} & \textbf{P(success ≥3)} & \textbf{P(Partial)} & \textbf{P(Miss)} & \textbf{Exp SB} \\
\midrule
10 & 99.8\% & <0.2\% & <0.1\% & 1.0 \\
\bottomrule
\end{tabular}
\end{table}
\textbf{Interpretation}: You will almost never Miss or Partial. Nearly every roll is Clean Success or Success with SB. Boons from failures vanish.

\subsection{SB Abundance}
With 10 dice, ~1 SB per roll. 4-player round → ~4 SB per round. GM spends aggressively.

\textbf{Heuristic}: Every success will likely have a cost. Expect every rolled action to generate a complication. Plan backup plans, escape routes, redundant assets.

\subsection{Boon Scarcity – Strategic Implications}
Boon sources: bonds (1–2 per session), GM awards. Miss rate near zero.

\textbf{Optimal Boon use}:
\begin{itemize}
    \item Never spend on re-rolling (you already succeed).
    \item Never convert to XP (waste).
    \item Spend to upgrade Position (Controlled → Dominant) to reduce expected 1s.
    \item Spend to activate Assets – primary use.
\end{itemize}

\subsection{Position – SB Mitigation}
Dominant reduces expected 1s by ~30\%. Desperate increases 1s → avoid.

\subsection{Clock Management}
Clocks size 8–10, multiple simultaneous. GM ticks clocks with SB. You cannot out-roll clocks. Use assets and narrative actions to manipulate them.

\subsection{Damage and Survival – Rocket Tag}
Tier IV foes: Harm 5–6, Armor 2. First strike wins. Build for assassination or unkillable defense.

\subsection{Skill Synergies – Diminishing Returns}
Beyond 10 dice, gains microscopic. Stop at Attribute 5 + Skill 5 = 10 dice. Add Cap 3 follower (+3 assist) = effective 13 dice (cap). Invest spare XP in talents, major assets, Cap 4–5 followers, prestige abilities.

\subsection{Assets and Followers – Force Multipliers}
At Tier III, have at least one Standard asset and one Cap 3–4 follower.

\subsection{Example Build Tier III (120 XP)}
From Tier II (80 XP), add 40 XP: Body 4→5 (15), Melee 4→5 (10), Battle Instincts (6), Cap 3 Follower (9) = 40 XP. Final dice: Body 5 + Melee 5 = 10 dice; with follower assist +3 = 13 effective.

\subsection{Tier III Heuristics Summary}
\begin{enumerate}
    \item Success is guaranteed. Question is ``at what cost?''
    \item SB abundant – expect complications on every roll.
    \item Boons extremely scarce – hoard for defensive use (upgrading Position).
    \item Dominant position mandatory – reduces expected 1s. Desperate is a trap.
    \item Clocks escalate rapidly – use assets to delay/reverse, not dice.
    \item Avoid damage, don't soak it – invest in defensive talents and surprise attacks.
    \item Dice past 10 give diminishing returns – stop at 10. Invest in +Effect talents and narrative assets.
    \item Followers are cheap dice – Cap 3 (9 XP) gives +3 assist better than Attribute 4→5 (15 XP for +1 die).
    \item Assets manipulate clocks – Standard asset (8 XP) can negate a complication per session.
    \item Game shifts from ``can I succeed?'' to ``how much will this cost me?'' Embrace the cost – it's the source of drama.
\end{enumerate}

\section{Tier IV (151–220 XP) Optimal Heuristics}

\subsection{Probability Landscape}
\begin{table}[h]
\centering
\begin{tabular}{lcccc}
\toprule
\textbf{Dice} & \textbf{P(success ≥3)} & \textbf{P(Partial)} & \textbf{P(Miss)} & \textbf{Exp SB} \\
\midrule
12 & ~99.98\% & <0.02\% & ~0.0001\% & 1.2 \\
\bottomrule
\end{tabular}
\end{table}
Success absolute. No Partials, no Misses. Boons only from bonds and GM awards (3–6 per session).

\subsection{SB Abundance – GM Toolkit}
12 dice → 1.2 SB per roll → ~5 SB per round. GM spends 2–4 SB per complication, often pooling SB from multiple rolls into major twists (4+ SB events).

\subsection{Boon Scarcity – Strategic Implications}
Treat each Boon as a minor miracle. Primary use: upgrade Position to reduce SB generation, or activate assets.

\subsection{Position – The Only Remaining Lever}
Dominant reduces expected 1s from 1.2 to ~0.8 (0.4 fewer SB per roll). Significant over session. Always seek Dominant.

\subsection{Clock Management – Multiple, Interlocking, Rapid}
Three or four clocks active simultaneously, size 8–10. GM ticks multiple each round. You cannot out-roll them. Use assets to tick clocks backward, sacrifice minor assets to pause clocks, split party to manage different clocks.

\subsection{Damage and Survival – Rocket Tag}
Tier IV foes: Harm 5–6, Armor 2, multiple attacks. Do not engage in fair fights. First strike wins. Build for assassination or unkillable defense.

Defensive talents: ``Unbreakable'' (ignore all Harm from one attack), ``Shadow Form'' (intangible), ``Final Witness'' (redirect killing blow). Offensive: ``Deathblow'' (triple Harm from stealth), ``Absolute Judgment'' (+3 dice), ``Market Tyrant'' (cripple logistics).

\subsection{Skill Synergies – Diminishing Returns}
Beyond 12 dice, gains microscopic. Stop at Attribute 5 + Skill 5 = 10 dice. Add Cap 3 follower = effective 13 dice. Spare XP (80–150 XP beyond Tier III) into multiple talents, major assets, Cap 4–5 followers, prestige abilities.

\subsection{Assets and Followers – True Power}
At Tier IV, have at least two Standard assets, one Major asset, and one Cap 4–5 follower. Your character sheet reflects sphere of influence, not personal stats.

\subsection{Example Build – Tier IV Commander (200 XP)}
From Tier III (120 XP), add 80 XP: Wits 3→4 (12), Presence 2→4 (21), talents (Inspiring Leader 8, Tactical Genius 10), Major Asset ``Fortified Keep'' (12), Cap 4 Follower ``Elite Squad'' (16) = 79 XP. Final dice: Body 5 + Melee 5 = 10 dice; squad assist +3 = 13 effective. Command: Presence 3 + Command 4 = 7 dice + squad assist = 10.

\subsection{Tier IV Heuristics Summary}
\begin{enumerate}
    \item Success absolute. Game now about managing SB and complications.
    \item Boons extremely scarce (3–6 per session). Use only for Position upgrade or asset activation.
    \item SB abundant (5+ per round). Expect every action to generate a moderate complication. Plan for cascading failures.
    \item Dominant position mandatory – reduces expected 1s by ~30\%. Desperate is death.
    \item Clocks rapid and multiple – cannot out-roll. Use assets and narrative actions.
    \item Combat is rocket tag – first strike wins. Build for assassination or unkillable defense.
    \item Diminishing returns past 12 dice – stop at Attribute 5 + Skill 5 + follower assist. Invest in talents and assets.
    \item Assets are primary tool – solve problems without rolling. Major asset worth more than +2 dice.
    \item Followers provide cheap dice – Cap 4 (16 XP) gives +3 assist, better than raising Attribute 4→5 (15 XP for +1 die).
    \item Embrace the cost – every victory leaves scars, debts, and new enemies. That's the point.
\end{enumerate}

% ========== NEW HEURISTICS SECTIONS ==========

\subsection*{Magic-Specific Heuristics (All Tiers)}

\textbf{Tier I:}
\begin{itemize}
    \item Stick to 1–2 TAGS per spell (DV 2–3). \texttt{[HEAL]} cantrip is your most efficient action (removes 1 Fatigue, no backlash risk).
    \item Save multi-tag spells for desperate moments.
    \item For Free Casters: Accept minor backlash (Fatigue 1) rather than avoid casting.
\end{itemize}

\textbf{Tier II:}
\begin{itemize}
    \item Dangerous tags (Leap, Fold, Gate, Transform, Dominate) cost +2 DV. Only use when you have 2+ Boons to improve Position or reroll.
    \item Acceptable backlash: Fatigue 1. Unacceptable: Harm or Corruption.
    \item Runekeepers: Clear Obligation after every session (Act of Service). Never let Obligation exceed Spirit+Presence – Fatigue penalties kill your Position.
\end{itemize}

\textbf{Tier III:}
\begin{itemize}
    \item Your dice pool (8–10) makes DV 5 spells reliable. But every 1 generates SB. Cast \texttt{[HEAL]} on allies \textit{before} they take Harm – Fatigue removal is cheap; Harm healing is expensive.
    \item Cantors: Corruption blooms are not failures – they're power spikes. Plan to bloom once per arc. The burden is a story hook, not a punishment.
\end{itemize}

\subsection*{Social Encounter Heuristics}

\begin{itemize}
    \item \textbf{Before rolling}, ask: ``What does this NPC want that they haven't said?'' A Wits+Insight roll (DV 3) reveals one hidden need. Use it – it can lower DV by 1 for future rolls.
    \item \textbf{Strings are social Boons.} Spend a String to auto-succeed on a minor social request (e.g., ``let me through,'' ``tell me a rumor''). Hoard 1–2 Strings for critical moments.
    \item \textbf{Offer a Debt Clock instead of a roll.} When an NPC refuses, say: ``I'll owe you.'' Start a 4-segment Debt Clock. When it fills, you must perform a service – but you get what you want now.
    \item \textbf{Tier III+ social:} You almost never fail. Instead, the GM spends SB to add complications (a witness overhears, a rival is alerted). Accept these – they're the price of success.
\end{itemize}

\subsection*{Travel and Exploration Heuristics}

\begin{itemize}
    \item \textbf{Travel roles:} Always assign a Guide (Wits+Survival). On a success, tick the Travel Clock +1. On a Partial, tick +1 but gain 1 Boon. On a Miss, tick +1 and the GM gains 1 SB.
    \item \textbf{Use followers for independent actions.} A Cap 3 scout can ``Scout \& Signal'' – change the party's next travel Position to Dominant. Cost: mark Exposure +1 on the follower.
    \item \textbf{Weather is Position.} In heavy rain or fog, default travel Position is Desperate. Spend 1 Boon to improve to Controlled (you find a sheltered path).
    \item \textbf{Tier III travel:} You will succeed. The question is how many SB the GM banks. Avoid rolling unless necessary – use Assets (maps, guides) to bypass rolls entirely.
\end{itemize}

\subsection*{Asset \& Follower Status Heuristics}
\begin{itemize}
  \item \textbf{Keep Assets at Stable or better.} Flourishing gives +1 die and +1 Yield; Stable is fine. Strained gives no Yield and –1 die – fix it within one downtime.
  \item \textbf{A Collapsed asset is a story hook, not a loss.} Restoring it requires a quest or major downtime. Use it to generate a session’s plot (e.g., “The safehouse was seized by the city guard – get it back.”)
  \item \textbf{Monitor Follower Loyalty before Fitness.} A Strained follower may still fight but can refuse orders; a Hurt follower (Fitness) is merely less effective. Prioritise loyalty missions (e.g., personal favours, public recognition) over healing.
  \item \textbf{Use Independent Actions sparingly – each one worsens Loyalty or Fitness.} Save Scout \& Signal for critical moments (e.g., before a boss fight). For routine scouting, use a Cap 3+ follower’s Off‑Screen Capability instead (costs 1 SB but no status loss).
  \item \textbf{Delegate Assets to a Cap 4–5 follower only when their Loyalty is Faithful.} If a follower’s Loyalty drops to Strained while managing delegated Assets, all those Assets immediately worsen one step. Check Loyalty before every downtime.
\end{itemize}

\subsection*{Player Trust Heuristics}
\begin{itemize}
  \item \textbf{Form a Trust at Tier II or III, not earlier.} Below Tier II, XP is too scarce; pooling 10+ XP delays personal advancement. At Tier III, shared assets become force multipliers.
  \item \textbf{Delegate upkeep to a Cap 5 follower as soon as possible.} A single Cap 5 follower can manage 4 Assets for free (using their own upkeep). This saves the party 4–12 XP per downtime.
  \item \textbf{Never put all Assets under one delegated follower.} If that follower’s Loyalty becomes Strained or Broken, all delegated Assets worsen immediately. Split delegation between two Cap 4 followers (2 Assets each) for redundancy.
  \item \textbf{Dissolve the Trust only if the campaign ends or a player leaves permanently.} A dissolved Trust returns no XP; instead, transfer Assets to individuals by paying half their XP cost (or accepting Obligation to a terrestrial patron).
  \item \textbf{Use the Trust pool for one‑time emergencies.} With majority agreement, spend 2 XP from the pool to activate a Trust Asset twice in a session (instead of once). This is cheaper than each player spending their own Boons.
\end{itemize}

\subsection*{Patron and Obligation Heuristics}

\begin{itemize}
    \item \textbf{Obligation Capacity = Spirit + Presence.} Never let total Obligation exceed this. Mark Fatigue when you do. If you exceed by 2+, trigger a Patron Intrusion – embrace it as a story hook.
    \item \textbf{Clear Obligation every downtime.} An Act of Service (described scene) clears 1 segment. Do this before accumulating debt.
    \item \textbf{Multiple Patrons (Invoker):} Each Symbol has its own Obligation track. Total Obligation is the sum. If you carry 3 Symbols, you'll hit capacity fast. Specialize in 1–2 Patrons.
    \item \textbf{When to accept a Patron Intrusion:} When the GM says ``your Patron demands a task,'' say yes. It's free adventure fuel. Refusing costs you 1 Fatigue and the Patron may withdraw a Rite until you comply.
\end{itemize}

\subsection*{Terrestrial Patron Heuristics}
\begin{itemize}
  \item \textbf{Accept favors only when the Debt Clock is low (0–2 segments).} A clock at 3+ segments is a ticking bomb – the patron will call in the debt soon.
  \item \textbf{Always clear at least 1 Obligation per downtime.} Terrestrial patrons are less patient than cosmic ones; skipping two cycles may trigger Sanction (e.g., revoked permits, bounty).
  \item \textbf{Use a String before a favor.} If you have a String (e.g., “I did the guild a favour last year”), spend it to reduce the Obligation cost of a new favor by 1 segment.
  \item \textbf{When a Debt Clock fills, offer a counter-proposal.} Patrons are rational. “I can’t assassinate the duke, but I can deliver proof of his embezzlement” may reframe the debt without penalty.
  \item \textbf{Sever ties only if you have a replacement patron.} Losing a Tier 2 patron without a backup leaves you vulnerable. Cultivate at least one other terrestrial patron before burning bridges.
\end{itemize}

\subsection*{Healing and Recovery Heuristics}

\begin{itemize}
    \item \textbf{Use \texttt{[HEAL]} after every fight.} One action, DV 2, removes 1 Fatigue. No backlash risk (unless you roll 1s, which give the GM SB – acceptable).
    \item \textbf{Accept Harm to reset Fatigue (the roller-coaster).} If you're at Fatigue 3 (Desperate, –2 dice), a small Harm 1 (e.g., trap, punch) clears Fatigue and leaves you at Harm 1, Fatigue 0. This is a tactical reset – use it.
    \item \textbf{Healing Harm is expensive.} DV 5+ for \texttt{[HEAL]+[PURIFY]}. Only do this if the target is at Harm 2 or dying. Otherwise, rest for a week (downtime).
    \item \textbf{Tier III+ healing:} You will succeed. The GM will spend SB on complications (scars, Patron attention). Accept these – they're better than a dead ally.
\end{itemize}

\subsection*{Combat Positioning and Movement Heuristics}

\begin{itemize}
    \item \textbf{Sidestep Strike:} Move one band as part of a melee attack, but Position worsens by one step. Use this to close from Near to Close with a light weapon (+2d) – the extra dice offset the Position penalty.
    \item \textbf{Charge (Full turn):} Move two bands and attack, cost 1 Fatigue, next defense Desperate. Only charge when you have a Boon to improve defense Position, or when you're sure to kill the target.
    \item \textbf{Opportunity attacks:} Leaving Close range without Disengage provokes a free melee attack. Always Disengage (action) unless you have a talent that ignores this (e.g., Tactical Movement).
    \item \textbf{Flanking:} If two allies are in Close with the same enemy, both gain +1 die to melee attacks. Coordinate.
\end{itemize}

\subsection*{Party Composition Synergy Heuristics}

\begin{itemize}
    \item \textbf{Essential roles (Tier I–II):} At least one \texttt{[HEAL]} user (removes Fatigue), one scout (Wits+Stealth), one face (Presence+Sway). Combat can be anyone.
    \item \textbf{Essential roles (Tier III+):} At least one Asset manager (keeps safehouse/spy network), one follower commander (Cap 3+ to delegate), one magic user (to handle wards/curses).
    \item \textbf{Skill overlap:} Two characters with Stealth 3+ is better than one with Stealth 5 – you can split the party for infiltration.
    \item \textbf{Bonds:} Ensure every PC has at least one Bond with another PC. That Bond grants 1 Boon per session when invoked. Use it.
\end{itemize}

\subsection*{GM-Facing Heuristics (Optional)}

\begin{itemize}
    \item \textbf{SB spending per scene:} Target 1–2 SB for Tier I, 2–4 for Tier II, 4–6 for Tier III+. Unused SB are wasted pressure.
    \item \textbf{Clock size by tier:} Tier I: 4–6 segments. Tier II: 6–8. Tier III: 8–10. Larger clocks = more guaranteed escalation.
    \item \textbf{When to worsen Position:} Do it when the fiction demands it (ambush, bad weather, social faux pas). Do not do it just to raise difficulty – use SB for that.
    \item \textbf{Deck of Consequences:} Draw 1 card per 2 SB spent (rounded up). Synthesize suits into a single complication. Do not draw for every 1 – it slows play.
\end{itemize}

\subsection*{Transitioning Between Tiers}

\begin{table}[h]
\centering
\begin{tabular}{ll}
\toprule
\textbf{From → To} & \textbf{Change in Mindset} \\
\midrule
Tier I → II & Stop hoarding Boons. You'll get fewer Misses. Spend Boons on Position upgrades before rolls. \\
Tier II → III & Success is nearly guaranteed. Accept that every roll will generate SB. Plan for complications, not failure. \\
Tier III → IV & Your personal dice are maxed. Invest in followers and assets – they're now your primary leverage. \\
Tier IV → V & You are a faction, not a person. Solve problems with assets and Debt Clocks, not dice. \\
\bottomrule
\end{tabular}
\end{table}

\section*{Closing Note}

These heuristics are grounded in the game's probability math and resource economy. They are not prescriptive rules but strategic observations. The best Fate's Edge players internalize these patterns and then break them for the sake of story. When reviewing the system, assess how well it enables this kind of informed, narrative-driven decision-making.

\end{document}